\documentclass[10pt,a4paper]{amsart}
\usepackage{amsmath,amssymb,amsthm,graphicx}
\usepackage{verbatim}
\newenvironment{reference}{\comment}{\endcomment}
%\newenvironment{reference}{}{}
\newenvironment{slogan}{\comment}{\endcomment}
\newenvironment{history}{\comment}{\endcomment}

% For commutative diagrams we use Xy-pic
\usepackage[all]{xy}

% We use 2cell for 2-commutative diagrams.
\xyoption{2cell}
% Theorem environments.
%
\theoremstyle{plain}
\newtheorem{theorem}[subsection]{Theorem}
\newtheorem{proposition}[subsection]{Proposition}
\newtheorem{lemma}[subsection]{Lemma}

\theoremstyle{definition}
\newtheorem{definition}[subsection]{Definition}
\newtheorem{example}[subsection]{Example}
\newtheorem{exercise}[subsection]{Exercise}
\newtheorem{situation}[subsection]{Situation}

\theoremstyle{remark}
\newtheorem{remark}[subsection]{Remark}
\newtheorem{remarks}[subsection]{Remarks}

\numberwithin{equation}{subsection}

% Macros
%
\def\lim{\mathop{\mathrm{lim}}\nolimits}
\def\colim{\mathop{\mathrm{colim}}\nolimits}
\def\Spec{\mathop{\mathrm{Spec}}}
\def\Hom{\mathop{\mathrm{Hom}}\nolimits}
\def\Ext{\mathop{\mathrm{Ext}}\nolimits}
\def\SheafHom{\mathop{\mathcal{H}\!\mathit{om}}\nolimits}
\def\SheafExt{\mathop{\mathcal{E}\!\mathit{xt}}\nolimits}
\def\Sch{\mathit{Sch}}
\def\Mor{\mathop{\mathrm{Mor}}\nolimits}
\def\Ob{\mathop{\mathrm{Ob}}\nolimits}
\def\Sh{\mathop{\mathit{Sh}}\nolimits}
\def\NL{\mathop{N\!L}\nolimits}
\def\CH{\mathop{\mathrm{CH}}\nolimits}
\def\proetale{{pro\text{-}\acute{e}tale}}
\def\etale{{\acute{e}tale}}
\def\QCoh{\mathit{QCoh}}
\def\Ker{\mathop{\mathrm{Ker}}}
\def\Im{\mathop{\mathrm{Im}}}
\def\Coker{\mathop{\mathrm{Coker}}}
\def\Coim{\mathop{\mathrm{Coim}}}

% Boxtimes
%
\DeclareMathSymbol{\boxtimes}{\mathbin}{AMSa}{"02}

%
% Macros for moduli stacks/spaces
%
\def\QCohstack{\mathcal{QC}\!\mathit{oh}}
\def\Cohstack{\mathcal{C}\!\mathit{oh}}
\def\Spacesstack{\mathcal{S}\!\mathit{paces}}
\def\Quotfunctor{\mathrm{Quot}}
\def\Hilbfunctor{\mathrm{Hilb}}
\def\Curvesstack{\mathcal{C}\!\mathit{urves}}
\def\Polarizedstack{\mathcal{P}\!\mathit{olarized}}
\def\Complexesstack{\mathcal{C}\!\mathit{omplexes}}
% \Pic is the operator that assigns to X its picard group, usage \Pic(X)
% \Picardstack_{X/B} denotes the Picard stack of X over B
% \Picardfunctor_{X/B} denotes the Picard functor of X over B
\def\Pic{\mathop{\mathrm{Pic}}\nolimits}
\def\Picardstack{\mathcal{P}\!\mathit{ic}}
\def\Picardfunctor{\mathrm{Pic}}
\def\Deformationcategory{\mathcal{D}\!\mathit{ef}}
\usepackage{cleveref}
\crefname{originalreference}{original reference}{original references}
\Crefname{originalreference}{Original reference}{Original references}
\makeatletter
\newlabel{definition-split-node}{{0CBV}{source}{}{sourcepointer-1}{}}
\newlabel{definition-split-node@cref}{{[originalreference][0][]0CBV}{[1][0][]source}{}{}{}}
\makeatother
\crefname{originallemma}{original lemma}{original lemmas}
\Crefname{originallemma}{Original lemma}{Original lemmas}
\makeatletter
\newlabel{divisors-lemma-base-change-and-fitting-ideal-omega}{{0C3I}{source}{}{sourcepointer-2}{}}
\newlabel{divisors-lemma-base-change-and-fitting-ideal-omega@cref}{{[originallemma][0][]0C3I}{[1][0][]source}{}{}{}}
\makeatother
\crefname{originalreference}{original reference}{original references}
\Crefname{originalreference}{Original reference}{Original references}
\makeatletter
\newlabel{lemma-descend-nodal-family}{{0C5F}{source}{}{sourcepointer-3}{}}
\newlabel{lemma-descend-nodal-family@cref}{{[originalreference][0][]0C5F}{[1][0][]source}{}{}{}}
\makeatother
\crefname{originalreference}{original reference}{original references}
\Crefname{originalreference}{Original reference}{Original references}
\makeatletter
\newlabel{lemma-formal-local-structure-nodal-family}{{0CBX}{source}{}{sourcepointer-4}{}}
\newlabel{lemma-formal-local-structure-nodal-family@cref}{{[originalreference][0][]0CBX}{[1][0][]source}{}{}{}}
\makeatother
\crefname{originalreference}{original reference}{original references}
\Crefname{originalreference}{Original reference}{Original references}
\makeatletter
\newlabel{lemma-nodal}{{0C4D}{source}{}{sourcepointer-5}{}}
\newlabel{lemma-nodal@cref}{{[originalreference][0][]0C4D}{[1][0][]source}{}{}{}}
\makeatother
\crefname{originallemma}{original lemma}{original lemmas}
\Crefname{originallemma}{Original lemma}{Original lemmas}
\makeatletter
\newlabel{limits-lemma-descend-finite-presentation}{{01ZM}{source}{}{sourcepointer-6}{}}
\newlabel{limits-lemma-descend-finite-presentation@cref}{{[originallemma][0][]01ZM}{[1][0][]source}{}{}{}}
\makeatother
\crefname{originalproposition}{original proposition}{original propositions}
\Crefname{originalproposition}{Original proposition}{Original propositions}
\makeatletter
\newlabel{more-algebra-proposition-ubiquity-G-ring}{{07PX}{source}{}{sourcepointer-7}{}}
\newlabel{more-algebra-proposition-ubiquity-G-ring@cref}{{[originalproposition][0][]07PX}{[1][0][]source}{}{}{}}
\makeatother
\crefname{originallemma}{original lemma}{original lemmas}
\Crefname{originallemma}{Original lemma}{Original lemmas}
\makeatletter
\newlabel{more-morphisms-lemma-realize-prescribed-residue-field-extension-etale}{{02LF}{source}{}{sourcepointer-8}{}}
\newlabel{more-morphisms-lemma-realize-prescribed-residue-field-extension-etale@cref}{{[originallemma][0][]02LF}{[1][0][]source}{}{}{}}
\makeatother
\crefname{originallemma}{original lemma}{original lemmas}
\Crefname{originallemma}{Original lemma}{Original lemmas}
\makeatletter
\newlabel{more-morphisms-lemma-relative-map-approximation-pre}{{0GDX}{source}{}{sourcepointer-9}{}}
\newlabel{more-morphisms-lemma-relative-map-approximation-pre@cref}{{[originallemma][0][]0GDX}{[1][0][]source}{}{}{}}
\makeatother
\crefname{originalreference}{original reference}{original references}
\Crefname{originalreference}{Original reference}{Original references}
\makeatletter
\newlabel{remark-quadratic-extension}{{0CBT}{source}{}{sourcepointer-10}{}}
\newlabel{remark-quadratic-extension@cref}{{[originalreference][0][]0CBT}{[1][0][]source}{}{}{}}
\makeatother
\title{Etale local structure of a family of curves at a singular point of the fibre: etale pullback to the standard uv - a family over Z[a]}
\author{The Stacks Project Authors}
\date{stacks-project:curves.tex:local-snapshot-20261012; GNU FDL 1.2 or later}
\begin{document}
\maketitle
\noindent\textit{Editorial note.} All delivered mathematical text is the author's own, copied from the pinned Stacks Project chapter file curves.tex; no statement or proof step was rewritten, reordered or removed.\par\medskip
\noindent\textit{Editorial note.} The counted proof is the author's complete argument: the approximation reduction with its base-change compatibility of Sing(f), the etale neighbourhood realizing kappa'/kappa(s) in both the split and field cases, the completed normal form st - h with the explicit fibre product presentation of Y\_sigma, and the closing Artin approximation step are all retained verbatim.\par\medskip
\noindent\textit{Editorial note.} The label lemma-etale-local-structure-nodal-family names the counted statement, and the proof environment immediately following it in the source is delivered verbatim as its proof.\par\medskip
\noindent\textit{Editorial note.} The counted statement is carried by the author's own lemma environment in the chapter source.\par\medskip
\noindent\textit{Allowed standard prerequisites (editorial).}\begin{itemize}
\item Absolute Noetherian approximation. Limits, Lemma limits-lemma-descend-finite-presentation, with Lemma lemma-descend-nodal-family and Divisors, Lemma divisors-lemma-base-change-and-fitting-ideal-omega. Supplies the reduction to S of finite type over Z: write R = colim R\_i, descend f to f\_i at-worst-nodal of relative dimension 1 with x\_i a singular point of the fibre.
\item Artin approximation for the completed local algebra. More on Morphisms, Lemma more-morphisms-lemma-relative-map-approximation-pre, with More on Algebra, Proposition more-algebra-proposition-ubiquity-G-ring. Turns the completed isomorphism O\_\{X,x\}\textasciicircum{}wedge = O\_\{Y\_sigma,y\_sigma\}\textasciicircum{}wedge into an etale-local isomorphism after an etale neighbourhood, since O\_\{S,s\} is a G-ring.
\end{itemize}\medskip
\ifdefined\hypertarget\hypertarget{sourcepointer-1}{}\fi
\noindent\textit{Source reference.} Original reference 0CBV: sources/stacks-curves/pinned-original-source.tex, source lines 4572--4572. Prerequisite pointer invoked at proof line(s) 5356; it names a standard fact of the same chapter that the counted argument uses, and every step of the counted proof built on it is retained verbatim in the delivered body.\par\medskip
\ifdefined\hypertarget\hypertarget{sourcepointer-2}{}\fi
\noindent\textit{Source reference.} Original lemma 0C3I: sources/stacks-chow/tags.txt, source lines 15193--15193. Prerequisite pointer invoked at proof line(s) 5329; it names a standard fact of a sibling chapter that the counted argument uses, and every step of the counted proof built on it is retained verbatim in the delivered body.\par\medskip
\ifdefined\hypertarget\hypertarget{sourcepointer-3}{}\fi
\noindent\textit{Source reference.} Original reference 0C5F: sources/stacks-curves/pinned-original-source.tex, source lines 5096--5096. Prerequisite pointer invoked at proof line(s) 5325; it names a standard fact of the same chapter that the counted argument uses, and every step of the counted proof built on it is retained verbatim in the delivered body.\par\medskip
\ifdefined\hypertarget\hypertarget{sourcepointer-4}{}\fi
\noindent\textit{Source reference.} Original reference 0CBX: sources/stacks-curves/pinned-original-source.tex, source lines 5137--5137. Prerequisite pointer invoked at proof line(s) 5357; it names a standard fact of the same chapter that the counted argument uses, and every step of the counted proof built on it is retained verbatim in the delivered body.\par\medskip
\ifdefined\hypertarget\hypertarget{sourcepointer-5}{}\fi
\noindent\textit{Source reference.} Original reference 0C4D: sources/stacks-curves/pinned-original-source.tex, source lines 4391--4391. Prerequisite pointer invoked at proof line(s) 5340; it names a standard fact of the same chapter that the counted argument uses, and every step of the counted proof built on it is retained verbatim in the delivered body.\par\medskip
\ifdefined\hypertarget\hypertarget{sourcepointer-6}{}\fi
\noindent\textit{Source reference.} Original lemma 01ZM: sources/stacks-chow/tags.txt, source lines 2599--2599. Prerequisite pointer invoked at proof line(s) 5321; it names a standard fact of a sibling chapter that the counted argument uses, and every step of the counted proof built on it is retained verbatim in the delivered body.\par\medskip
\ifdefined\hypertarget\hypertarget{sourcepointer-7}{}\fi
\noindent\textit{Source reference.} Original proposition 07PX: sources/stacks-chow/tags.txt, source lines 9812--9812. Prerequisite pointer invoked at proof line(s) 5384; it names a standard fact of a sibling chapter that the counted argument uses, and every step of the counted proof built on it is retained verbatim in the delivered body.\par\medskip
\ifdefined\hypertarget\hypertarget{sourcepointer-8}{}\fi
\noindent\textit{Source reference.} Original lemma 02LF: sources/stacks-chow/tags.txt, source lines 3406--3406. Prerequisite pointer invoked at proof line(s) 5345; it names a standard fact of a sibling chapter that the counted argument uses, and every step of the counted proof built on it is retained verbatim in the delivered body.\par\medskip
\ifdefined\hypertarget\hypertarget{sourcepointer-9}{}\fi
\noindent\textit{Source reference.} Original lemma 0GDX: sources/stacks-chow/tags.txt, source lines 20457--20457. Prerequisite pointer invoked at proof line(s) 5386; it names a standard fact of a sibling chapter that the counted argument uses, and every step of the counted proof built on it is retained verbatim in the delivered body.\par\medskip
\ifdefined\hypertarget\hypertarget{sourcepointer-10}{}\fi
\noindent\textit{Source reference.} Original reference 0CBT: sources/stacks-curves/pinned-original-source.tex, source lines 4521--4521. Prerequisite pointer invoked at proof line(s) 5343; it names a standard fact of the same chapter that the counted argument uses, and every step of the counted proof built on it is retained verbatim in the delivered body.\par\medskip
\begin{lemma}
\label{lemma-etale-local-structure-nodal-family}
Let $f : X \to S$ be a morphism of schemes. Assume that
$f$ is at-worst-nodal of relative dimension $1$. Let
$x \in X$ be a point which is a singular point of the
fibre $X_s$. Then there exists a commutative diagram of schemes
$$
\xymatrix{
X \ar[d] &
U \ar[rr] \ar[l] \ar[rd] & &
W \ar[r] \ar[ld] &
\Spec(\mathbf{Z}[u, v, a]/(uv - a)) \ar[d] \\
S & &
V \ar[ll] \ar[rr] & & \Spec(\mathbf{Z}[a])
}
$$
with $X \leftarrow U$, $S \leftarrow V$, and $U \to W$ \'etale morphisms,
and with the right hand square cartesian, such that there exists
a point $u \in U$ mapping to $x$ in $X$.
\end{lemma}

\begin{proof}[First proof using Artin approximation]
We first use absolute Noetherian approximation to reduce to the
case of schemes of finite type over $\mathbf{Z}$.
The question is local on $X$ and $S$. Hence we may assume that
$X$ and $S$ are affine. Then we can write $S = \Spec(R)$
and write $R$ as a filtered colimit $R = \colim R_i$
of finite type $\mathbf{Z}$-algebras.
Using Limits, Lemma \ref{limits-lemma-descend-finite-presentation}
we can find an $i$ and a morphism $f_i : X_i \to \Spec(R_i)$ whose base
change to $S$ is $f$. After increasing $i$ we may assume that $f_i$
is at-worst-nodal of relative dimension $1$, see
Lemma \ref{lemma-descend-nodal-family}.
The image $x_i \in X_i$ of $x$ will be a singular
point of its fibre, for example because the formation of
$\text{Sing}(f)$ commutes with base change (Divisors, Lemma
\ref{divisors-lemma-base-change-and-fitting-ideal-omega}).
If we can prove the lemma for $f_i : X_i \to S_i$ and
$x_i$, then the lemma follows for $f : X \to S$ by base
change. Thus we reduce to the case studied in the next
paragraph.

\medskip\noindent
Assume $S$ is of finite type over $\mathbf{Z}$. Let $s \in S$ be the
image of $x$. Recall that $\kappa(x)$ is a finite separable extension
of $\kappa(s)$, for example because $\text{Sing}(f) \to S$
is unramified or because $x$ is a node of the fibre $X_s$
and we can apply Lemma \ref{lemma-nodal}.
Furthermore, let $\kappa'/\kappa(x)$ be the
degree $2$ separable algebra associated to $\mathcal{O}_{X_s, x}$ in
Remark \ref{remark-quadratic-extension}.
By More on Morphisms, Lemma
\ref{more-morphisms-lemma-realize-prescribed-residue-field-extension-etale}
we can choose an \'etale neighbourhood $(V, v) \to (S, s)$
such that the extension $\kappa(v)/\kappa(s)$ realizes either
the extension $\kappa(x)/\kappa(s)$ in case
$\kappa' \cong \kappa(x) \times \kappa(x)$ or
the extension $\kappa'/\kappa(s)$ if $\kappa'$ is a field.
After replacing $X$ by $X \times_S V$ and $S$ by $V$
we reduce to the situation described in the next paragraph.

\medskip\noindent
Assume $S$ is of finite type over $\mathbf{Z}$ and
$x \in X_s$ is a split node, see Definition \ref{definition-split-node}.
By Lemma \ref{lemma-formal-local-structure-nodal-family} we see that there
exists an $\mathcal{O}_{S, s}$-algebra isomorphism
$$
\mathcal{O}_{X, x}^\wedge \cong
\mathcal{O}_{S, s}^\wedge[[s, t]]/(st - h)
$$
for some $h \in \mathfrak m_s^\wedge \subset \mathcal{O}_{S, s}^\wedge$.
In other words, if we consider the homomorphism
$$
\sigma : \mathbf{Z}[a] \longrightarrow \mathcal{O}_{S, s}^\wedge
$$
sending $a$ to $h$, then there exists an $\mathcal{O}_{S, s}$-algebra
isomorphism
$$
\mathcal{O}_{X, x}^\wedge
\longrightarrow
\mathcal{O}_{Y_\sigma, y_\sigma}^\wedge
$$
where
$$
Y_\sigma = \Spec(\mathbf{Z}[u, v, t]/(uv - a))
\times_{\Spec(\mathbf{Z}[a]), \sigma} \Spec(\mathcal{O}_{S, s}^\wedge)
$$
and $y_\sigma$ is the point of $Y_\sigma$ lying over
the closed point of $\Spec(\mathcal{O}_{S, s}^\wedge)$
and having coordinates $u, v$ equal to zero. Since
$\mathcal{O}_{S, s}$ is a G-ring by
More on Algebra, Proposition \ref{more-algebra-proposition-ubiquity-G-ring}
we may apply More on Morphisms, Lemma
\ref{more-morphisms-lemma-relative-map-approximation-pre}
to conclude.
\end{proof}

\end{document}
